\chapter{Constitutive Formulations \Paper{4/5}}

\section{Introduction}

Recall that up until this point we have regarded the beam as a purely one-dimensional body: the state of stress $\FrameResult(\xi_m)$ and strain $\FrameStrain(\xi_m)$ at a point $\xi_m$ are mathematically related only to the material points immediately to their left $\xi_m-d\xi$ and right $\xi_m+d\xi$.

However, in reality, a beam always has \emph{some} material emanating from the centerline $\bm{x}$, and ideally the dependence of $\FrameResult$ on $\FrameStrain$ should take into account the distribution of this
additional material. To this end, we formally introduce the material points $\hat{\bm{\rho}}$ in the cross section through an extended three-dimensional displacement field $\hat{\FramePosition}(\reflenx, \hat{\bm{\rho}})$ which is assumed to depend only on the configuration variables of the reference curve \(\chi=\bigl(\FramePosition(\reflenx), \FrameRotation (\reflenx), \bm{\alpha}(\reflenx)\bigr)\) at the 1D material point $\reflenx$. Because the beam strains \(\FrameStrain(\reflenx) = (\ShearStrain , \CurveStrain , \bm{\alpha}^{\prime})\) have been formulated as unique invariant functions of the configuration variables $\chi$, it follows that an invariant strain measure for the extended deformation $\hat{\bm{x}}(\chi)$ should be expressible as functions of only the beam strains $\FrameStrain$.

\begin{itemize}
\item All basic equations are formulated with respect to an arbitrary cartesian coordinate system which is not restricted to principal axes. Thus the origin of this system is not necessarily a special point like the centroid. This relieves the input of the finite element data.
\end{itemize}

\section{Kinematic Framework}

The extended displacement field for the newly introduced section material points $\hat{\rho}$  is given by:
\begin{equation*}
\hat{\bm{x}}(\xi, \hat{\rho}) 
= \bm{x} + \FrameRotation \bigl(\PlanePosition (\hat{\rho}) + (\bm{\alpha}\cdot\bm{\varphi})\, \mathbf{D}\bigr)
\end{equation*}
%
where the 2D vector $\PlanePosition  = y\mathbf{D}_{2} + z\mathbf{D}_3$ identifies a section material point $\hat{\rho}$.

\subsection{Deformation}

\(\hat{\FramePosition} = \hat{\chi}(\xi,\rho)\), consider the deformation gradient \(\bm{F} \triangleq T\hat{\chi}\).

\subsection{Constitution}
The second Piola-Kirchhoff stress \(\bm{S}\)


\subsection{Conjugation}
Recall the Frobenius inner product for two second order tensors \(\bm{A}\) and \(\bm{B}\) with coordinate representations \(A^{ij} \triangleq \mathbf{G}^i\cdot \bm{A}\mathbf{G}^j\) and  \(B_{ij} \triangleq \mathbf{G}_i\cdot \bm{B}\mathbf{G}_j\)
\begin{equation*}
    \left<\bm{A},\,\bm{B}\right> = \operatorname{tr} \bm{A}\bm{B}^{\mathrm{t}} = A^{ij} B_{ij}
\end{equation*}
If \(\delta \SolidStrain = \left(\hat{\bm{B}} \delta\FrameStrain\right) \stackrel{s}{\otimes}\mathbf{D}\) for some second-order tensor \(\hat{\bm{B}}\), then
\begin{equation*}
    \left<\bm{S},\,\delta\SolidStrain\right> = \SectionIntegral{ 
    \operatorname{tr} \left[\bm{S}\frac{1}{2}\left(\hat{\bm{B}}\delta \FrameStrain\otimes\mathbf{D} + \mathbf{D}\otimes \hat{\bm{B}}\delta \FrameStrain\right)
    \right]
    }
\end{equation*}
Exploiting linearity of the trace together with the identity \(\operatorname{tr} (\bm{a}\otimes\bm{b}) = \bm{a}\cdot \bm{b}\) for arbitrary vectors \(\bm{a}\) and \(\bm{b}\) it follows
\begin{equation*}
    \left<\bm{S},\,\delta\SolidStrain\right> = \delta \FrameStrain\cdot \SectionIntegral{
        \hat{\bm{B}}^{\mathrm{t}}\frac{1}{2}\left(\bm{S}\mathbf{D} + \bm{S}^{\mathrm{t}}\mathbf{D}\right)
    }
\end{equation*}
and consequently by requiring \(\left<\FrameResult,\,\delta\FrameStrain\right>=\left<\bm{S},\,\delta\SolidStrain\right>\)
\begin{equation*}
\FrameResult \equiv \SectionIntegral{\hat{\bm{B}}^{\mathrm{t}}\frac{1}{2}\left(\bm{S}\mathbf{D} + \bm{S}^{\mathrm{t}}\mathbf{D}\right)}
\end{equation*}
For symmetric \(\bm{S}\) one further has
\begin{equation}
\FrameResult \equiv \SectionIntegral{
  \hat{\bm{B}}^{\mathrm{t}} \bm{S}\mathbf{D}
}
\end{equation}
This has the intuitive implication that ony the first ``column'' and/or ``row'' of the stress tensor \(\bm{S}\) should contribute to the frame stress resultants \(\ShearResult\) and \(\CurveResult\).

\subsection{Constraints}
Following the treatment by \cite{steigmann2017notes} who treated incompressibility. 
\emph{Three} constraints :
\begin{equation}
    \mathbf{D}_{\alpha} \cdot \bm{S}\mathbf{D}_{\beta} = 0
\end{equation}
Differentiating
\begin{equation*}
    \frac{\mathrm{d}}{\mathrm{d} \tau} \mathbf{D}_{\alpha} \cdot \bm{S}\mathbf{D}_{\beta} 
    =  \mathbf{D}_{\alpha} \cdot \mathbb{C}\dot{\SolidStrain}\mathbf{D}_{\beta} = 0
\end{equation*}

\subsection{Linearization}
\begin{equation*}
    D {\FrameResult} [\delta \FrameStrain] = \hat{\bm{B}}^{\mathrm{t}}D \bm{S} [D \SolidStrain [\delta \FrameStrain]]
\end{equation*}
Using \(\mathbb{C}[\delta \SolidStrain] \triangleq D\bm{S}[\delta \SolidStrain]\) for vectors \(\delta \SolidStrain\) tangent at \(\SolidStrain\) 
\begin{equation*}
 D {\FrameResult} [\delta \FrameStrain] = \hat{\bm{B}}^{\mathrm{t}} \mathbb{C} D \SolidStrain [\delta \FrameStrain]
\end{equation*}
Now using \(D \SolidStrain [\delta \FrameStrain] = (\hat{\bm{B}}\delta \FrameStrain) \stackrel{s}{\otimes} \mathbf{D}\)
\begin{equation*}
\hat{\mathbf{K}} = \SectionIntegral{
    \hat{\mathbf{B}}^{\mathrm{t}} \mathbb{C}\hat{\bm{B}} + \hat{\bm{G}}
    }
\end{equation*}

\subsection{Summary}
\begin{ambox}[Cauchy Section Constitution]{box:cauchy-section}

\begin{itemize}
    \item Strain tensor \(\SolidStrain = \SolidStrain(\ShearStrain,\CurveStrain,\WarpStrain,\WarpAmplitude)\)
    \item Cauchy constitution \(\bm{S} = \bm{S}(\SolidStrain)\)
    \item The conjugate stress resultant \(\FrameResult\) is obtained by integration over the cross section
    \[
    \FrameResult= \SectionIntegral{\hat{\bm{B}}^{\mathrm{t}} \bm{S}\mathbf{D}}\]
    where \(\hat{\bm{B}}\) is defined by the directional derivative \(D\SolidStrain[\delta\FrameStrain] = \left(\hat{\bm{B}}\delta\FrameStrain\right)\stackrel{s}{\otimes}\mathbf{D}\)
\end{itemize}
\end{ambox}

\section{Hierarchical Kinematics}

It will be convenient to introduce the relative displacement $\bm{\pi}$ given by the expression:
$$
\bm{\pi} \triangleq \FrameRotation^{\mathrm{t}} (\hat{\bm{x}} - \bm{x}) = \PlanePosition  + (\bm{\alpha}\cdot\bm{\varphi})\,\mathbf{D}
$$

\subsection{Exact Kinematics}

The \emph{exact} deformation gradient for the extended material field $\hat{\bm{x}}$ is:
$$
\bm{F}(\xi, \hat{\rho}) = \nabla \hat{\bm{x}} = \FrameRotation\left[\bm{1} + \bm{\varepsilon}\otimes\mathbf{D}\right] + \FrameRotation\nabla \bm{\omega}
$$
where
$$
\bm{\varepsilon}(\hat{\rho}) \triangleq\ShearStrain  + \CurveStrain  \times (\PlanePosition  + \bm{\omega})
\qquad\text{ and }\qquad
\bm{\omega} \triangleq (\bm{\alpha}\cdot\bm{\varphi})\, \mathbf{D}
$$
and $\ShearStrain $ and $\CurveStrain $ are standard beam strains.
Note that in the presence of warping $\bm{\varepsilon}$ is nonlinear in
the strain variables $\FrameStrain$ and may be written as the sum of linear $\bar{\bm{\varepsilon}}$ and quadratic $\CurveStrain \times\bm{\omega}$ terms:
$$
\bm{\varepsilon} = \bar{\bm{\varepsilon}} + \CurveStrain \times\bm{\omega}
\qquad\text{ with }\qquad
\bar{\bm{\varepsilon}} \triangleq \ShearStrain  + \CurveStrain \times\PlanePosition 
$$
%
The exact Green-Lagrange strain tensor may be expressed as:
\begin{equation}
\boxed{
\SolidStrain = \bm{\varepsilon}
\stackrel{\mathrm{s}}{\otimes}
\mathbf{D} 
+\nabla\omega \stackrel{\mathrm{s}}{\otimes} \mathbf{D}
+({\bm{\varepsilon}}\cdot\mathbf{D})\,\nabla\omega \stackrel{\mathrm{s}}{\otimes} \mathbf{D}
+ \frac{1}{2}\left(
  \bm{\varepsilon}\cdot\bm{\varepsilon}\,\mathbf{D}\otimes\mathbf{D}
+ \nabla\omega\otimes\nabla\omega
\right)
}
\end{equation}
%
Some important, progressively simpler, special cases include:
\begin{enumerate}
  \item \textbf{Torsion Warping}: In this case only a single warping mode is considered so that \(\bm{\alpha}\) is a scalar and the warping mode \(\varphi\) is St. Venant's warping function:
  $$
  \nabla \omega = \bm{\Phi}\begin{pmatrix}\alpha^{\prime} \\ \alpha \end{pmatrix}
  \qquad\text{ with }\qquad
  \bm{\Phi} \triangleq \begin{bmatrix}
  \varphi & 0 \\ 0 & \varphi_{, y} \\ 0 & \varphi_{, z}
  \end{bmatrix}
  $$
  In the absence of constraints on \(\alpha\), this corresponds to the RBV formulation described by \cite{armero2021modeling}.
  This model is well suited for implementation in \gls{frame:shear} models. Examples of finite element implementations include \cite{simo1991geometricallyexact} and \cite{gruttmann2000theory}.

  \item \textbf{Constrained Warping} Referred to as the TWKV formulation by \cite{armero2021modeling}, this formulation imposes the constraint 
  \[
  \alpha = \mathbf{D} \cdot \CurveStrain
  \]
  This assumption is well suited from implementation with \gls{frame:euler} models. Examples include \cite{nukala2004mixed}, \cite[Section 3.1.2]{lecorvec2012nonlinear} and \cite{zang2011formulation} for symmetric sections, whose work is implemented in OpenSees and derives from \cite{alemdar2001distributed}.

\item \textbf{Uniform torsion}: In this case warping is assumed to depend on the torsional component of the curvature $\CurveStrain $ so that $\alpha = \mathbf{D}\cdot\CurveStrain $ and $\alpha^{\prime}=0$. Therefore:
  $$
  \bm{\omega} = (\varphi \mathbf{D}\otimes\mathbf{D})\CurveStrain  = \tau \varphi \mathbf{D}
  $$
  and
  $$
  \bm{\varepsilon} = \ShearStrain  + \CurveStrain \times\PlanePosition  + \tau\varphi \CurveStrain \times\mathbf{D}
  $$

\item \textbf{Plane Section}: In the absence of warping one has $\bm{\omega}=\bm{0}$ and $\bm{\varepsilon}\equiv \bar{\bm{\varepsilon}}$:
  $$
  \SolidStrain = \bar{\bm{\varepsilon}}
  \stackrel{\mathrm{s}}{\otimes}
  \mathbf{D} 
  + \frac{1}{2}\left.
    \bar{\bm{\varepsilon}}\cdot\bar{\bm{\varepsilon}}\,\mathbf{D}\otimes\mathbf{D}\right. \\
  $$
\end{enumerate}


\subsection{Linear Kinematics}
Even for geometrically exact models, it is common to assume linear section kinematics (see, e.g., \cite{simo1991geometricallyexact}).
% \subsection{Shear Correction}

$$
\begin{aligned}
\hat{\bm{B}}  = \left[\bm{\Psi} ~~~~ [\PlanePosition ^{\times}]^{\mathrm{t}} ~~~~ \bm{\Phi}\right]
= \left[\begin{array}{ccc|ccc|cc}
& \ShearStrain  & & & \CurveStrain  & & \alpha^{\prime} & \alpha \\ \hline
1 &                0  &                0 &  0 & z  & -y & \tilde{\varphi}(y, z) &                                 \\
0 & \psi_{y y}(y, z)  & \psi_{y z}(y, z) & -z & 0  &  0 &            0 & \tilde{\varphi}_{, y}(y, z) \\
0 & \psi_{z y}(y, z)  & \psi_{z z}(y, z) &  y & 0  &  0 &            0 & \tilde{\varphi}_{, z}(y, z) \\
\end{array}\right]
\end{aligned}
$$

\section{Resultants}

\subsection{Solid}
\subsubsection{Linear}
\begin{equation*}
= \SectionIntegral{
\begin{bmatrix}
\bm{\Psi}^{\mathrm{t}}\mathbb{C}\bm{\Psi} & -\bm{\Psi}^{\mathrm{t}}\mathbb{C}[\bm{r}^\times] & \bm{\Psi}^{\mathrm{t}}\mathbb{C}\bm{\Phi} \\
. & -\left[\bm{r}^{\times}\right]\mathbb{C}[\bm{r}^\times] & \left[\bm{r}^{\times}\right]\mathbb{C}\bm{\Phi} \\
. & . & \bm{\Phi}^{\mathrm{t}}\mathbb{C}\bm{\Phi} \\
\end{bmatrix} 
}
\end{equation*}

\subsection{Isotropy}
\subsubsection{Linear}

\begin{equation*}
=\begin{bmatrix}
\bm{\Psi}^{\mathrm{t}}\bm{\Psi} & -\bm{\Psi}^{\mathrm{t}}[\bm{r}^\times] & \bm{\Psi}^{\mathrm{t}}\bm{\Phi} \\
. & -\left[\bm{r}^{\times}\right][\bm{r}^\times] & \left[\bm{r}^{\times}\right]\bm{\Phi} \\
. & . & \bm{\Phi}^{\mathrm{t}}\bm{\Phi} \\
\end{bmatrix} 
\end{equation*}

The final stiffness matrix for a general isotropic cross section is
\begin{equation}
\hat{\mathbf{K}} = \left[\begin{array}{ccc|ccc|cc}
E A & 0 & 0 & 0 & E Q_{y} & E Q_{z} & E R_{w} & 0\\
0 & G A_{y} & 0 & G Q_{yx} & 0 & 0 & 0 & G R_{y}\\
0 & 0 & G A_{z} & G Q_{zx} & 0 & 0 & 0 & G R_{z}\\ \hline
0 & G Q_{yx} & G Q_{zx} & G I_{0} & 0 & 0 & 0 & G S_{a} \\
E Q_{y} & 0 & 0 & 0 & E I_{y} & - E I_{yz} & E S_{y} & 0 \\
E Q_{z} & 0 & 0 & 0 & - E I_{yz} & E I_{z} & E S_{z} & 0 \\ \hline
E R_{w} & 0 & 0 & 0 & E S_{y} & E S_{z} & E C_{w} & 0 \\
0 & G R_{y} & G R_{z} & G S_{a} & 0 & 0 & 0 & G C_{a}
\end{array}\right]
\end{equation}

For the case constrained by \(\alpha = \mathbf{D}\cdot \CurveStrain\) one has:
\begin{equation*}
\hat{\mathbf{K}} = \left[\begin{array}{ccc|ccc|c}
E A & 0 & 0 & 0 & E Q_{y} & E Q_{z} & E R_{w}\\
0 & G A_{y} & 0 & G Q_{yx} + G R_{y} & 0 & 0 & 0\\
0 & 0 & G A_{z} & G Q_{zx} + G R_{z} & 0 & 0 & 0\\ \hline
0 & G Q_{yx} + G R_{y} & G Q_{zx} + G R_{z} & G (C_{a} + I_{0} + 2 S_{a}) & 0 & 0 & 0\\
E Q_{y} & 0 & 0 & 0 & E I_{y} & - E I_{yz} & E S_{y}\\
E Q_{z} & 0 & 0 & 0 & - E I_{yz} & E I_{z} & E S_{z}\\ \hline
E R_{w} & 0 & 0 & 0 & E S_{y} & E S_{z} & E C_{w}
\end{array}\right]
\end{equation*}
Note that this generalizes the work of \cite{saritas2006mixed}, as presented in Equation~3.23 of \cite{dire20173d}.

\section{Summary}

\section{Examples}


\begin{exmpl}[End-loaded cantilever with an angle (asymmetric) cross section]{ex:angle-section}
\end{exmpl}


\noindent %
\begin{exmpl}[Modeling of a tapered bridge column]{ex:const-flare}
\end{exmpl}

\input{Part I/Sections/A_Flexure}

\section{St. Venant Torsional Warping}\label{sec:torsion-warping}

The objective of this section is to detail the computation of St. Venant's warping function \(\varphi\) and related cross section parameters \(C_{\varphi}, C_w, J, ...\).

\subsection{The Neumann Problem}
Consider a body with reference configuration 
\begin{equation*}
    \bm{x}_0(\reflenx) + \SectionCoordinate(\varsigma)
\end{equation*}
undergoing a deformation with the form
\begin{equation*}
    \bm{x}(\reflenx, \varsigma ) = \bm{x}_0(\reflenx) + \SectionCoordinate(\varsigma) + \vartheta \, \FrameBasis\times(\SectionCoordinate - \bm{\SectionLocation}) + \alpha \varphi(\varsigma) \FrameBasis
\end{equation*}
Subject to the traction-free boundary condition
\begin{equation*}
\begin{aligned}
    \bm{S}\mathbf{n} &= \bm{0}
\end{aligned}
\end{equation*}
where \(\mathbf{n}\) is the unit normal on the boundary in the section's plane. 
The displacement field \(\bm{u}\) is 
\begin{equation*}
\bm{u}(\reflenx, y, z) =
\vartheta \, \FrameBasis\times(\SectionCoordinate - \bm{\SectionLocation}) + \alpha \varphi(\SectionCoordinate) \FrameBasis
\end{equation*}
and the infinitesimal strain vector
\begin{equation*}
\bm{\varepsilon} = \vartheta' \FrameBasis \times (\SectionCoordinate - \bm{\SectionLocation}) + \alpha' \varphi \FrameBasis + \alpha \nabla \varphi
\end{equation*}
The traction-free boundary condition implies
\begin{equation*}
    \nabla_{\mathbf{n}}\bm{\varepsilon} = \bm{0}
\end{equation*}
Recall \cref{eq:saint-venant-stress} and that for a prismatic section the outward normal is orthogonal to the axial basis \(\FrameBasis\), \(\mathbf{n}\cdot \FrameBasis = 0\).
\begin{marsdenbox}
\noindent
Torsional warping problem for rotation about some point \(\bm{\SectionLocation}\)
\begin{equation}
\begin{array}{rll}
\Delta \varphi_{\SectionLocation} &=0 & \text { in } \Omega, \\
\nabla_{\mathbf{n}} \varphi_{\SectionLocation}  &=-\alpha \FrameBasis \times (\mathbf{r}-\bm{\SectionLocation}) \cdot \mathbf{n} & \text { on } \partial \Omega  \\
\SectionIntegral{\varphi_{\SectionLocation}} &= 0
\end{array}
\label{eq:torsion-neumann}
\end{equation}
where $\mathbf{n}$ is the unit outward normal to the boundary of the cross section domain.
\end{marsdenbox}


\subsection{Center of Twist}
A natural choice for \(\bm{\SectionLocation}\) is to select it such that the following conditions hold:
\begin{equation}
    \SectionIntegral{\varphi_{\SectionLocation} (\bm{r}-\bm{\SectionLocation})\times\FrameBasis} = \bm{0}
    \label{eq:shear-center-constraint}
\end{equation}
Thus, the point \(\tilde{\bm{\SectionLocation}}\) is defined as that for which the corresponding solution \(\tilde{\varphi}\) to \cref{eq:torsion-neumann} with \(\pi = \tilde{\bm{\SectionLocation}}\) satisfies \cref{eq:shear-center-constraint}. 
Note that \cref{eq:shear-center-constraint} simply requires that the \(m-w\) coupling tensor of \cref{eq:def-linear-isotropic-coupling-tensors} vanishes.
...
say
\[
\tilde{\varphi} \triangleq \varphi - [\tilde{\bm{\SectionLocation}} , \FrameBasis, (\bm{r} - \hat{\bm{\SectionLocation}})]
\]
In components 
\[
\tilde{\varphi} = \varphi - 
\]
or equivalently
\[
\tilde{\varphi}\bm{r}\times\FrameBasis = \varphi\bm{r}\times\FrameBasis - \bm{r}\times\FrameBasis \, \left(\tilde{\bm{\SectionLocation}} \cdot \FrameBasis \times \hat{\bm{r}}\right)
\]
\[
\tilde{\varphi}\bm{r}\times\FrameBasis = \varphi\bm{r}\times\FrameBasis - \bm{r}\times\FrameBasis \otimes \left(\FrameBasis \times \hat{\bm{r}}\right) \tilde{\bm{\SectionLocation}}
\]
Using this in \cref{eq:shear-center-constraint} furnishes the linear system:
\begin{equation}
    \SectionIntegral{\tilde{\varphi}\,\bm{r}\times\FrameBasis}
    =\left(\SectionIntegral{\bm{r}\times\FrameBasis\otimes\FrameBasis\times\bar{\bm{r}}} \right)\tilde{\bm{\SectionLocation}}
    \label{eq:derive-shear-center-1}
\end{equation}
where \(\bar{\bm{r}} \triangleq \bm{r} - \bar{\bm{\SectionLocation}}\).
Also recall that 
\[
\SectionIntegral{\varphi_{\SectionLocation} \bm{r}\times\FrameBasis} 
\equiv \SectionIntegral{\varphi_{\SectionLocation} (\bm{r}-\bar{\bm{\SectionLocation}})\times\FrameBasis} 
%\equiv \bm{0}
\]
Thus we use \(\bar{\bm{r}}\) for \(\bm{r}\) in \cref{eq:derive-shear-center-1}
\begin{equation}
    \SectionIntegral{\tilde{\varphi}\,\bar{\bm{r}}\times\FrameBasis}
    =\left(\SectionIntegral{\bar{\bm{r}}\times\FrameBasis\otimes\FrameBasis\times\bar{\bm{r}}} \right)\tilde{\bm{\SectionLocation}}
\end{equation}
Pulling out the constant \(\FrameBasis\) 
\begin{equation}
    \SectionIntegral{\tilde{\varphi}\,\bar{\bm{r}}\times\FrameBasis}
    =\FrameBasis^{\times}\left(\SectionIntegral{\bar{\bm{r}}\otimes\bar{\bm{r}}} \right)\FrameBasis^\times\tilde{\bm{\SectionLocation}}
\end{equation}


\subsection{Finite Elements}\label{finite-elements}

\cite{surana1979isoparametric}

\begin{quote}
\[
\begin{gathered}
\bm{K}_{a b}=\int \left(\frac{\partial N_{i}}{\partial x} \frac{\partial N_{b}}{\partial x}+\frac{\partial N_{i}}{\partial y} \frac{\partial N_{b}}{\partial y}\right) \operatorname{det}[J] \mathrm{d} \xi \mathrm{d} \eta \\
\bm{f}_{a}=\SectionIntegral{\left(y \frac{\partial N_{i}}{\partial x}-x \frac{\partial N_{i}}{\partial y}\right) \operatorname{det}[J]
} \\
\left\{\delta_{a}\right\}=\left\{\psi_{a}\right\}
\end{gathered}
\]
\end{quote}


\begin{equation}
    \bm{f}\cdot\bm{w} = \FrameBasis \cdot \SectionIntegral{(\bm{B}\bm{w} \times \FrameBasis_{\beta} \bm{N}\bm{x}^e_{\beta})}
\end{equation}

\begin{equation}
    \bm{f}\cdot\bm{w} = -\FrameBasis \cdot \FrameBasis_{\beta}\times \SectionIntegral{ (\bm{N}\bm{x}^e_{\beta})\, \bm{B}\bm{w}}
\end{equation}